\section{语言模型里的 RL 设定}

\subsection{state, action, reward}

把整段 response 记作一个动作 $a$，语言模型 RL 的对应关系非常直接：
\begin{itemize}
    \item \textbf{State $s$}：prompt 加上已经生成的 response 前缀。
    \item \textbf{Action $a$}：生成下一个 token，或者在 outcome reward 视角里，直接把整段 response 看成一个动作序列。
    \item \textbf{Reward $R(s,a)$}：整段 response 的好坏，本文只讨论 \textbf{outcome reward} 和 \textbf{verifiable reward}。
\end{itemize}

字幕里强调了一点：在语言模型场景中，transition dynamics 基本就是拼接操作。你做出的下一状态就是
\[
s' = s + a
\]
也就是把 token 继续 append 到上下文里。这和机器人控制非常不同，后者的状态由物理世界决定，语言模型的状态则更像是“自己写出来的中间草稿”。

\begin{knowledgebox}{语言模型里的 state 很“虚”}
在 robotics 中，state 可能是关节角、速度、图像；在 LM 中，state 只是 token 序列。这个差别意味着：
\begin{itemize}
    \item LM 可以自由构造 scratchpad；
    \item LM 可以在 test time 做 planning / search；
    \item 但 reward 往往更稀疏，正确与错误的边界也更硬。
\end{itemize}
\end{knowledgebox}


\begin{warningbox}{process reward 不是白送的}
如果能在中间步骤给出可信的 reward，训练通常会更好；但在语言模型里，“中间步骤是否对”往往比最终答案更难判断。因此本讲先聚焦 outcome reward，把机制讲清楚。
\end{warningbox}

\subsection{为什么 outcome reward 很自然}

这一讲选择 outcome reward，不只是因为它简单，而是因为它和课程脚本里的 toy task 完全一致：你给定一个 prompt，模型生成一个完整 response，最后用一个确定性的函数判断它是否“做对了”。

这类 reward 的优点是实现成本低，且容易复现：
\begin{itemize}
    \item 你可以直接写一个 Python 函数来算分；
    \item 训练和评估使用同一个标准；
    \item 不需要额外训练 reward model；
    \item 很适合数学、代码、排序、格式控制这类任务。
\end{itemize}

\begin{importantbox}{课程里为什么先讲这种 reward}
因为它最能暴露 policy gradient 的本质问题：reward 稀疏、方差大、更新不稳定。只要把这类 task 讲明白，后面更复杂的 RLHF / RLAIF 也更容易理解。
\end{importantbox}

\begin{table}[H]
\centering
\begin{tabular}{@{}p{3.1cm}p{4.4cm}p{5.0cm}@{}}
\toprule
\textbf{reward 类型} & \textbf{优点} & \textbf{缺点} \\
\midrule
outcome / verifiable reward & 可自动计算，标准明确 & 反馈稀疏，难学 \\
process reward & 中间步骤也能给分 & 需要判断过程本身是否可信 \\
learned reward model & 能吸收人类偏好 & 额外系统复杂度高 \\
\bottomrule
\end{tabular}
\caption{为什么课程这里先采用 outcome reward：它最利于把机制讲透}
\end{table}

\subsection{目标函数}

如果把 prompt 视作由环境给出的 $s$，response 视作 policy 采样出来的 $a$，那么目标就是最大化期望奖励：
\[
\max_\pi \; \mathbb{E}_{s \sim p(s),\, a \sim \pi(\cdot \mid s)}[R(s,a)].
\]

在 outcome reward 设定下，reward 只在整段输出结束后给出，因此 discounting、bootstrapping 这些经典 RL 概念在这里不那么核心。

\begin{knowledgebox}{为什么这里几乎不谈 discount}
如果一个 response 的好坏只能在最后统一打分，那么 reward 的主要难点已经不是时间折扣，而是稀疏性和方差。也就是说，\emph{什么时候} 给分不如 \emph{怎么把这一个分数变得可学} 更重要。
\end{knowledgebox}


